% REV_P0 2026-11-24：由 chapter_03_methods_rewrite_v1_20261122.tex 派生；V1保持不变。
% 新增：与原作者DLPFC基准的预处理/K/重复结构逐项对照；B15D提名经过；AutoK重复次数差异；
% BayesSpace与GMM K-blind规则缺陷；技术seed SD的含义。未改动任何冻结协议参数或实验数据。
% 第三篇CellCharter毕业论文第三章（2026-11-22）新版首稿
% 研究对象：DLPFC 12切片 + Wu等外部6患者，不涉及本账号其他论文。
% 保留 chapter_03_methods_integrated.tex 历史原稿；本稿不覆盖旧稿。
% 本文件是待嵌入学校主LaTeX模板的正文片段；依赖 booktabs、amsmath 等宏包。
% 真实结果数值见 chapter_04_results_rewrite_v1_20261122.tex 和来源审查矩阵。
% E3d后验ARI/NMI与E3b容差是模型已完成、研究者已看过其他结果后进行的事后分析。
\section{研究方法与实验方案}
\label{sec:cellcharter-methods-v1}

\subsection{研究思路与技术路线}
\label{subsec:cc-overview-v1}

空间转录组测序技术可以在组织位置上同时获得基因表达信息，为空间域（spatial domain）识别提供了数据基础。空间聚类既受到原始表达特征的影响，也受到邻接结构、聚类数选择和组织标注粒度的影响。因此，仅报告单一算法在某张切片上的最高ARI不足以反映实际结果的稳定性。本研究以成熟的CellCharter邻域聚合与聚类实现为基础，重点分析不同尺度表示、聚类数确定方式和评价定义对空间域聚类结论的影响；同时设置非空间及空间成熟方法作为对照，并使用不同组织来源的数据进行独立队列评估。

研究程序分为六个相互衔接但证据性质不同的阶段：（1）在12张DLPFC切片上构建并审计5维scVI潜在表示及25维CellCharter多尺度缓存；（2）固定Oracle-K与10技术种子比较10种邻域特征组织方式（E1/E2）；（3）以预定K网格检查尺度结论对簇数的敏感性（E3c），再以无标签AutoK实现独立选择K（E3d）；（4）从已保存的预测与原始六邻接阵列坐标重新计算严格空间边界和容差指标（E3b）；（5）在相同DLPFC切片上对照四种成熟基线（E5a）；（6）在6名乳腺癌患者的独立来源数据上重新执行同类方法，报告跨组织实验流程的表现（E5b）。E4空间图及标记基因表达图用于辅助解释，不参与前述特征或K参数的选择。

以上各阶段不是将一个在DLPFC训练完成的聚类模型迁移至乳腺癌，也不构成对CellCharter底层邻域聚合公式的原创改造。各阶段的原始冻结协议、预测及指标文件、计算脚本和审计记录均以项目实际归档版本为准。对实验结束后才新增的诊断性评价，本文在方法和结果中单独说明为事后分析，不将其倒写为早期已预注册的主研究终点。

\subsection{数据来源、样本层级与参考标注}
\label{subsec:cc-data-v1}

\subsubsection{DLPFC人脑组织切片}
DLPFC评价集由12张Visium空间转录组切片组成，切片标识依次为151507、151508、151509、151510、151669、151670、151671、151672、151673、151674、151675、151676。样本按来源分为3位供体，每位供体4张切片。既有归档供体映射中，151507--151510对应一组，151669--151672对应一组，151673--151676对应另一组；部分早期切片的供体对应关系依据公开附录佐证而非本研究重新取得并验证的全部原始受试者元数据，因此在最终统计解释中保留该证据限制。

每张切片采用原始10x \texttt{filtered\_feature\_bc\_matrix.h5}表达计数文件，空间位置信息来自\texttt{spots\_aligned.csv}，其中的barcode、\texttt{array\_row}、\texttt{array\_col}及像素坐标与源数据按唯一barcode匹配。人工参考组织层标签从\texttt{manual\_layer\_annotations.csv}读取，并根据\texttt{sample\_name}与\texttt{spot\_name}进行对齐。标注缺失的spot可保留于模型输入与空间可视化，但不纳入ARI、NMI或有效参考边界的计分。不得把未经重新病理判读的公开参考标签表述为本研究新制作的组织学金标准。

151507、151508的早期开发性参考文件与后续10片的标准化标签来源路径不完全相同。后续10张切片涉及已有\texttt{layer\_guess}及其顺序调整字段的核对；本研究按冻结的标准化标签与评价mask复用，不因结果重新修改。曾发现151507目录中存在与151669重复的错误\texttt{metadata.tsv}副本，按原数据审计记录该文件不被当前E1/E2聚类评价任务直接读取；这不等于证明所有更早的特征缓存生成过程完全不受任何数据卫生风险影响。上述问题作为来源质量与解释限制保留，不以人工修补方式改变本次已冻结预测。

\subsubsection{乳腺癌独立队列}
外部评价采用Wu等研究公开的6例人乳腺癌Visium数据，患者标识为1142243F、1160920F、CID4290、CID4465、CID44971及CID4535。源数据包含表达计数、空间坐标和供评价使用的组织区域标签；原始来源及压缩文件的MD5哈希已在实验前核对，患者标识互不相同。本研究的外部数据意在检验同类无监督聚类流程与邻域深度比较在另一组织来源下的表现，而非验证皮层空间域模型的直接跨器官预测能力，也不是医学诊断性能验证。

根据冻结输入审计记录，6名患者分别具有4784、4895、2426、1211、1160和1125个用于模型计算的spot；对应评价mask包含4665、4846、2419、1207、1159和1028个spot，合计15601个模型spot和15324个评价spot。评价spot数量不等于生物学重复次数。本文以6名不同患者作为独立生物学单位，spot、技术seed、尺度条件和预测任务均为患者内观察。

\subsubsection{样本独立性与数据使用边界}
DLPFC的独立生物学层级为供体$n=3$，每供体4张切片；外部数据的独立生物学层级为患者$n=6$。在计算时可以报告每切片或每患者的单独分数及其技术重复间变异，但不以同一供体的4张切片、同一患者的多个spot或同一实验条件的多个随机种子扩增统计检验样本量。DLPFC包含开发性探索过程，不将其12张切片全部描述为严格未接触的测试样本；外部队列在原方案冻结之后进行，但后续事后统计摘要仍须如实披露时间线。

\subsection{表达预处理与scVI特征缓存}
\label{subsec:cc-preprocess-v1}
DLPFC采用现有已审计的源处理和特征生成结果，按唯一barcode对齐空间位置表及表达矩阵。此前准备脚本将基因名规范为唯一，按计数阈值进行过滤并验证spot/barcode约束；在\texttt{counts}层保留原始计数，进行总量归一化与对数转换，并以基于原计数的\texttt{seurat\_v3}方法选择5000个高变基因。scVI通过\texttt{scvi.model.SCVI.setup\_anndata(..., layer="counts")}读取计数层，生成潜变量维度为5的表示。此处scVI拟合输入为计数层，而不是误将已归一化对数矩阵称为直接的原始计数输入。

DLPFC来源训练脚本的主要设置包括随机种子22911、最大400训练轮次和early stopping；新增10张切片源流程曾要求CUDA环境。训练轮次上限不意味着每张切片都运行满400轮。每张切片形成的低维潜变量矩阵记为
\begin{equation}
Z=[z_1^\top,\ldots,z_N^\top]^\top\in\mathbb{R}^{N\times5},
\end{equation}
其中$N$为相应切片纳入建模的spot数。既有5维与25维特征缓存均伴随spot顺序、源哈希和生成记录保存；后续E1/E2、E3c、E3d从同一冻结特征缓存抽取所需列，不因评价数值重新训练scVI。

上述DLPFC特征流程（\texttt{scripts/cellcharter/prepare\_dlpfc\_10\_features.py}）与CellCharter官方单切片基准脚本\texttt{src/benchmarking/CellCharter/individual.py}\cite{cellcharteranalyses2024}的主要步骤相同，包括：\texttt{filter\_genes}与\texttt{filter\_cells}的\texttt{min\_counts=3}、总量归一化至$10^6$并取对数、在计数层上以\texttt{seurat\_v3}选5000个高变基因、5维scVI潜变量、AdamW优化器配合early stopping与\texttt{reduce\_lr\_on\_plateau}、4层邻域聚合。表\ref{tab:cc-official-protocol-rev-p0}逐项列出两者已知的差异。由于存在这些差异，即使两者数值相近，也不能视为逐spot等价的复现。

\begin{table}[htbp]
\centering
\caption{本研究DLPFC 25D流程与CellCharter官方单切片基准的设置对照（依据双方脚本；未核对的项目注明“未核对”）}
\label{tab:cc-official-protocol-rev-p0}
\small
\begin{tabularx}{\linewidth}{p{2.6cm}XX}
\toprule
项目 & 官方单切片基准\cite{cellcharteranalyses2024} & 本研究E1/E2\\
\midrule
切片范围 & 12张均运行；原文图件只报告9张测试切片 & 12张\\
聚类方式 & 单切片（\texttt{individual.py}）；原文主结论为9切片联合聚类 & 单切片\\
K & 151669--151672为5，其余为7 & 同左（Oracle-K）\\
重复结构 & 10次，每次同时更换scVI与GMM的种子 & scVI固定种子22911训练一次；10个GMM种子\\
scVI轮次 & 默认上限，early stopping & 上限400，early stopping\\
GMM设置 & \texttt{cc.tl.SpatialCluster}默认参数，GPU & CellCharter GMM接口，CPU，容差$10^{-3}$，正则$10^{-6}$，最多100 epoch\\
输入文件 & HumanPilot预处理h5ad & 10x \texttt{filtered\_feature\_bc\_matrix.h5}与\texttt{spots\_aligned.csv}\\
参考标签字段 & \texttt{sce.layer\_guess} & \texttt{layer\_guess}经\texttt{manual\_layer\_annotations.csv}对齐\\
软件版本 & 未核对 & 记录于各缓存\texttt{config.json}\\
评价指标 & 仅ARI & ARI、NMI，另有事后边界指标\\
\bottomrule
\end{tabularx}
\end{table}

乳腺癌队列则在每名患者各自的计数数据上重新完成标准化输入与scVI五维表征，再建立患者内CellCharter空间邻域聚合特征。外部批次的模型训练、计算设备和输出状态以\texttt{e5b\_scvi\_feature\_batch\_v1.py}及相应Windows审计为准。该过程并未加载DLPFC训练的scVI权重做跨组织迁移；因此跨数据集比较评价的是流程的重复应用，而非既训练模型的domain generalization。

\subsection{空间六邻接图与CellCharter多尺度表示}
\label{subsec:cc-spatial-v1}

使用Visium阵列中spot的\texttt{array\_row}、\texttt{array\_col}建立空间邻接图$G=(V,E)$，其中每个节点对应一个spot。对有效直接邻居，阵列偏移满足
\begin{equation}
(|\Delta r|,|\Delta c|)\in\{(0,2),(1,1)\}.
\label{eq:cc-hex-v1}
\end{equation}
上述条件对应六角阵列的六种最近邻方向；边按无向关系记录，排除自环、重复坐标和不匹配barcode。该图按空间物理布局建立，不以表达量距离定义基因表达kNN图。不同方法如BANKSY和BayesSpace会按照各自原生实现使用空间坐标，其空间邻居定义不应与本节CellCharter的六邻接边混为同一算法。

本研究沿用CellCharter官方实现的\texttt{aggregate\_neighbors}操作，以5维scVI表征为输入，聚合4层并将原始层与聚合层拼接成
\begin{equation}
H=[L_0\mid L_1\mid L_2\mid L_3\mid L_4]
\in\mathbb{R}^{N\times25},\qquad
L_\ell\in\mathbb{R}^{N\times5}.
\label{eq:cc-H-v1}
\end{equation}
$L_0$为未聚合的5维输入层，$L_1$--$L_4$为CellCharter程序输出的相应邻域块。这里的层级是软件特征块索引，不将某一层直接赋予已验证的组织学影响半径；也不臆造CellCharter内部未逐项核实的加权公式。缓存中的列顺序、spot对齐和SHA由前期脚本及审计记录锁定。

固定的累计邻域特征为$L_0$（5D）、$[L_0,L_1]$（10D）、$[L_0,L_1,L_2]$（15D）、$[L_0,\ldots,L_3]$（20D）、$[L_0,\ldots,L_4]$（25D）。另设5种等维但不连续的15D组合，分别为$[L_0,L_1,L_3]$、$[L_0,L_1,L_4]$、$[L_0,L_2,L_3]$、$[L_0,L_2,L_4]$和$[L_0,L_3,L_4]$。其中标准15D配置A为$[L_0,L_1,L_2]$，候选B为$[L_0,L_1,L_4]$。A/B保持输入维度一致，但不代表其他未控制的模型差异全部消失；候选B亦不构成新的CellCharter核心模型。

需要说明两点数据使用情况。其一，4层邻域、5000高变基因、5维潜变量这组参数，是CellCharter原作者在151507、151672、151673上按参考标签ARI调出的\cite{varrone2024}；本研究直接沿用，因此这3张切片对25D配置不是未接触样本。其二，候选B是在开发切片151507（Oracle-K=7，GMM种子22911）上看到B相对A的ARI高出0.047后才被提名的；151507与151508都属于开发阶段已分析过的切片，随后又和其余10张切片一起纳入E1/E2完整评价。因此，E1/E2中B与A的比较带有提名切片的选择偏倚。结果部分同时报告预设的12切片汇总和排除这两张开发切片后的事后敏感性汇总，两者分别解释。

\subsection{DLPFC冻结聚类实验设计}
\label{subsec:cc-dlpfc-protocol-v1}

\subsubsection{E1/E2：十种特征配置与十技术种子}
E1/E2固定12张DLPFC切片、10种特征配置和10个随机种子，共形成$12\times10\times10=1200$个切片--配置--seed评价任务。预定种子序列为22911、7449、25842、10379、6690、26128、21059、22159、32389和12815。所有条件使用CellCharter的GMM聚类接口（CPU、full covariance、K-means初始化、收敛容差$10^{-3}$、协方差正则$10^{-6}$、最多100 epoch），统一记录拟合结果、实际簇数、评价mask和源文件哈希。为避免重计历史运行，若个别已存任务与冻结配置和文件SHA完全一致，则审计合格后复用；1200指完整任务网格而非全部1200项均在同一批次重新计算。

此阶段使用参考标签类别数确定K：151669--151672的4张切片设$K=5$，其余8张设$K=7$。这是明确的\textit{Oracle-K}条件。单个spot的层标注不用于GMM参数估计，但使用参考类别数本身已经带入了标签先验，不能称其为完全不依赖参考标签的聚类数选择。各模型先落盘预测及其SHA后再以固定标签mask进行指标评价，不依据ARI/NMI回选最有利种子。

\subsubsection{E3c：固定K敏感性网格}
E3c进一步在12张切片、4种累计尺度（10/15/20/25D）、5个指定K（5至9）和前述10种子上建立$12\times4\times5\times10=2400$格敏感性网格。其中480格与已通过审计的E1/E2 Oracle-K条件重合，予以复用；其余1920格为新形成的评价结果。固定K网格是对预设聚类数的受控比较，不意味着算法已自动估计未知K；报告每张切片、每个K和每个尺度的全部ARI、NMI及负结果，而不依据参考标签选择最终推荐K。E3c与E1/E2使用相同冻结特征缓存和对应GMM配置。

\subsubsection{E3d：无标签AutoK稳定性选K}
E3d独立使用12切片和四个累计尺度，构成48个K-blind AutoK配置。冻结候选集合为$K\in\{2,3,\ldots,10\}$，CellCharter实现中同时拟合辅助K=1与11；每个内部K进行5次GMM初始化，正式使用CPU、满5次稳定性尝试与明确seed22911--22915。按稳定性均值最大确定K，并保留原程序的首个最大值并列处理。每组合内部记录$11\times5=55$次模型拟合标签档案，全部48组合对应2640份内部运行标签阵列。该过程使用固定输入及无标签稳定性，人工组织层标签不进入拟合、早停或K选择。

E3d与原作者的稳定性选K示例有两点不同。第一，本研究每个候选K只重复5次，原作者在DLPFC上使用10次（原文Fig.~1a说明与Methods）；原作者还在Supplementary Information中把重复次数增加到25次，检查10次是否已使FMI曲线收敛。本研究的5次重复未做这一收敛检查，因此不能排除重复次数不足造成的选K不稳定。第二，原作者是对12张切片\textbf{联合}运行选K，本研究是对每张切片\textbf{单独}运行。E3d的结果应理解为“5次重复、逐切片”条件下的表现，不等同于原作者推荐设置下的表现。

AutoK输出在完成之后冻结：保存每个切片--尺度的原选中K、各候选K的稳定性、唯一barcode对应的聚类预测、版本与来源SHA；先进行不看参考标签的独立文件审计。直到48项原始预测全部审计通过，才单独使用固定参考标签mask对其计算ARI/NMI。这一指标计算是在观察过其他实验结果后追加的\textbf{事后后验评价}；不能倒称E3d原冻结训练阶段已经完成标签评分。对于选中K=2或10的边界情形，如实记录而不事后扩K或重新拟合。

\subsection{空间边界与拓扑容差的事后评价（E3b）}
\label{subsec:cc-boundary-v1}

\subsubsection{严格边界与评价mask}
从已冻结的E1/E2 1200份Oracle-K聚类预测和同源参考层标签出发，利用Visium六邻接图构建无向有效边集
\begin{equation}
E^\ast=\{(i,j)\in E\mid y_i\text{与}y_j\text{均有参考标注}\},
\end{equation}
其中$y_i$是spot $i$的参考层标签，$\hat{y}_i$为预测的原始聚类ID。对任意有效边$(i,j)$，若$y_i\ne y_j$则标记真实参考边界，若$\hat{y}_i\ne\hat{y}_j$则标记预测边界。两集合交集是严格真正例$TP_0$，预测集合减去交集为$FP_0$，参考集合减去交集为$FN_0$。严格边界F1定义为
\begin{equation}
F_{1,0}=\frac{2TP_0}{2TP_0+FP_0+FN_0}.
\label{eq:cc-boundary-f1-v1}
\end{equation}
Precision和recall分别由TP/$(TP+FP)$和TP/$(TP+FN)$计算，边界空集合/零分母按评价程序的固定规则处理并保留相应计数。每个spot与每条空间边只能按对应切片的有效参考mask计入，不因预测分布发生改变而增减困难spot。

\subsubsection{1-hop、2-hop容差与一对一匹配}
严格边界要求预测和参考边界出现在完全同一条网格边，可能对轻微位置偏移敏感。作为原始实验完成后明确追加的\textbf{事后敏感性分析}，本研究在只包含有效参考端点的边构成的线图（line graph）上，分别设置0、1、2跳邻近容差。若两条无向原图边在端点上相接，则在线图上相邻；两条边在线图中的最短距离作为容差距离。

对容差$t\in\{0,1,2\}$，在预测边界和参考边界之间仅保留线图距离不超过$t$的候选匹配，进行最大基数的一对一二部匹配，从而保证一条参考边界至多与一条预测边界对应。由匹配数$TP_t$以及全部预测/真实边界数得到$FP_t$与$FN_t$，进而计算$F_{1,t}$。本研究不将1-hop或2-hop视为严格空间定位，不用宽松容差代替0-hop主结果，也不在观察不同$t$分数之后选择最有利的$t$作为唯一报告。

完整E1/E2包括1200份源预测，三个容差各对应1200条指标，合计3600条原始评价记录。E3d另以48份K-blind AutoK预测执行相同边界定义，形成独立的144项补充结果；两组由于K策略不同不得合并为统一的独立样本。E3b的新增评价不包含重新训练GMM、改变原预测或重新制作参考层标签，所有原预测和参考文件与既有审计SHA对照；其结果属于事后尺度与边界灵敏度证据。

\subsection{DLPFC成熟方法基线（E5a）}
\label{subsec:cc-e5a-v1}

E5a在同一12张DLPFC切片上设置四种对照算法：非空间高斯混合模型GMM、非空间Leiden、BANKSY和BayesSpace。每方法、每切片、每K模式使用3个固定技术种子（22911、7449、25842），同时报告K-blind与Oracle-K两模式，共$12\times4\times2\times3=288$份正式预测。固定主对照特征为CellCharter A15D，B15D作为独立补充而非在看到基线结果后选择最佳变体。基线均采用同一切片的有效spot及相同标签mask评分，但并非全部共享输入特征或计算预算，因此其性能差异是不同原生方法流程的比较而非完全等价的单因素消融。

非空间GMM与Leiden使用已冻结scVI的原始$L_0$五维表示；GMM采用scikit-learn全协方差、K-means初始化、容差$10^{-3}$、协方差正则$10^{-6}$、\texttt{n\_init=1}、最多100次迭代。Leiden在15近邻表示图上聚类。BANKSY使用原始计数和spot像素空间坐标，原生固定15空间近邻、最高角度阶$m=1$、$\lambda=0.2$、20个主成分，经其自身原生变换和Leiden聚类。BayesSpace以同源Visium原计数和\texttt{array\_row}/\texttt{array\_col}阵列坐标为输入，使用2000高变基因、15个主成分、原生\texttt{spatialPreprocess}及\texttt{spatialCluster}流程，正式MCMC配置含10000次迭代、1000次burn-in、\texttt{model=t}与\texttt{init.method=mclust}。BayesSpace最初存在R CSC矩阵索引排序的导入兼容错误，模型拟合前即中止；之后仅修正输入稀疏矩阵表示并保留旧失败记录，未更改冻结科学方法参数。该BayesSpace配置与CellCharter原作者基准中的BayesSpace配置不同：原作者调参选定1000个高变基因、15个主成分、$\gamma=3$、5000次迭代\cite{varrone2024,cellcharteranalyses2024}；本研究为2000个高变基因、15个主成分、\texttt{model=t}、10000次迭代，$\gamma$取默认值。两组BayesSpace数值因此不能直接比较。此外，本研究对照集\textbf{不包含}原文中单切片表现最好的STAGATE\cite{dong2022stagate}，这是对照覆盖范围的局限。

本研究预先定义两类K规则，见表\ref{tab:cc-baseline-K-v1}。Oracle-K仅使用参考标签的类别数量作为目标，不允许据ARI优化分辨率；部分Leiden或BANKSY配置可能无法恰好得到目标簇数，需记录实际K。对于BayesSpace，自动K-blind使用候选q=2至10的\texttt{qTune}平均伪对数似然最大值规则。此为本研究冻结的自动判据，\textbf{与BayesSpace官方建议的elbow判断不一致}\cite{bayesspacevignette}。对于非空间GMM，K-blind使用BIC最小值规则。两种基于似然的规则在实际运行中都大量选中搜索上界$K=10$（结果见第\ref{subsec:cellcharter-e5a-v1}节及REV\_P0汇总表T5）。这些结果全部保留，并作为K-blind可比性的限制；不根据正式评价分数事后改变规则或搜索区间。

\begin{table}[htbp]
\centering
\caption{E5a/E5b成熟基线的聚类数选择策略（方法选择与标签评价分离）}
\label{tab:cc-baseline-K-v1}
\begin{tabular}{lll}
\toprule
方法 & K-blind & Oracle-K\\
\midrule
GMM & $K=2,\ldots,10$的BIC最小值 & 固定参考类别数\\
Leiden & 固定resolution=1.0 & 固定16档分辨率中实际K最接近目标\\
BANKSY & 固定resolution=1.0 & 固定16档分辨率中实际K最接近目标\\
BayesSpace & qTune最大伪对数似然（q=2--10） & 固定参考类别数q\\
\bottomrule
\end{tabular}
\end{table}

其中16档Leiden分辨率为0.1、0.15、0.2、0.25、0.3、0.4、0.5、0.65、0.8、1.0、1.25、1.5、2.0、2.5、3.0和4.0。基线输入、条码、K策略及各方法源代码SHA在288项独立只读审计中均作了验证。不同方法训练失败或输入兼容性错误均保留时间线，不在模型表现报告中删除不利结果。

\subsection{乳腺癌独立队列评价（E5b）}
\label{subsec:cc-e5b-v1}

E5b在所有原定模型首次拟合之前冻结了六患者名单、CellCharter特征组合、四成熟基线、K策略、种子及评价顺序。每名患者单独从其原始计数重新拟合scVI五维表示并构建CellCharter 25维聚合缓存。CellCharter采用$L_0$、10D、A15D、20D、25D和B15D共6种配置，按患者$\times$尺度执行无标签稳定性AutoK选择，再在K-blind、Oracle-K两臂及3技术seed下分别拟合GMM，形成$6\times6\times2\times3=216$份CellCharter预测。外部AutoK延用K2--10的冻结候选范围和5次稳定性技术重复，K在该患者/尺度中决定后才对全部3种正式聚类seed固定使用。人工参考区域的类别数只作为Oracle-K先验，不进入K-blind分支的模型训练和K选择。

四成熟基线的外部实现延续E5a所定义的原生表达预处理、空间表示、K-blind与Oracle-K原则，各方法$6\times2\times3=36$份预测，共144份。于是外部总任务数为216+144=360份。所有预测必须先写入包含完整模型spot的barcode和原簇编号的独立输出文件，审计通过后才能读取事前隔离的\texttt{reference\_only}区域标注及冻结评价mask计算ARI/NMI。独立只读审计总计核验360份预测与360组后验分数；保留原始源和预测SHA、K选择轨迹、失败记录与每患者三种子的技术变异。

外部评价先以同一患者内部的3个技术seed取均值，再以6名患者等权平均形成宏平均。相同患者上不同方法或尺度的差值采用配对描述，保留每患者原始方向和幅度；患者间差异不能用若干千spot作独立重复。本文同时报告全部K-blind方法配置，并将Oracle-K列为明确的标签类别数知情敏感性对照。不使用外部患者的评价分数改变先前已冻结尺度、K候选或方法清单。任何在读取结果以后另行加入的bootstrap、不确定性范围或对比优先级均明确属于\textbf{事后补充统计}，不得声称原冻结协议事前规定。

\subsection{评价指标、汇总方式与完整性控制}
\label{subsec:cc-metrics-v1}

采用ARI与算术平均NMI评价有参考标注的spot，评价时只以相同参考mask比较不同方法。两指标均对预测聚类ID的重命名不敏感，原始聚类编号与参考组织层并非天然一一对应。ARI通过配对样本划分的一致性进行机会校正；NMI衡量预测聚类与参考标签之间的共享信息程度。因两者不直接评价空间边界的几何位置，本研究同时提供严格边界与固定容差F1；其定义和一对一匹配规则已在前节给出。评价过程中不得按最终数值对预测簇进行有标签重排、丢弃不利spot或重训模型。

统计遵循\textbf{技术重复先汇总，独立生物来源后描述}原则。各结果表同时给出切片间标准差（或范围）、3位供体各自的均值，以及切片内技术种子标准差的平均值，这些都只作描述性不确定性展示，不用于显著性检验。DLPFC的技术种子标准差只反映GMM初始化的变化（scVI固定种子），因此小于原作者同时更换scVI与GMM种子时的波动（官方单切片基准平均约0.051，本研究25D约0.028），不能解读为本流程“更稳定”。所有新增汇总均由\texttt{REV\_P0\_20261124/rev\_p0\_derived\_summaries\_v1.py}从既有CSV只读计算，输入与输出SHA256记录在\texttt{derived/MANIFEST\_SHA256.csv}中。DLPFC E1/E2/E3c的10个seed只度量重复拟合变异，E5a和E5b采用3技术seed；AutoK内部5次拟合用于无标签稳定性选择，不是5份独立患者结果。DLPFC 12张切片嵌套于3名供体，切片等权平均是描述性摘要，独立供体间泛化不确定性不能由12切片的简单标准误替代。外部队列为6名患者，患者宏平均及各患者差值均需报告，尤其避免只呈现平均较高的单个患者或最有利的技术seed。本研究未据上述层级进行把单spot、空间边、seed或切片误当独立生物学重复的显著性检验；事后bootstrap若呈现，也必须注明探索性质与小样本解释限制。

为保证可复核性，各阶段独立记录原数据、输入表示、位置表、参考标注、K选择规则、随机种子、软件脚本SHA、完整预测和结果汇总。E1/E2完整1200份源预测已上传GitHub并保留各源SHA清单；E3c的1920新增分数及完整K网格汇总、E3d的48份AutoK原审计及事后ARI/NMI、E3b全量3600条边界指标、E5a的288条基线评分、E5b的360条外部评分均已有对应结果包。此事实表示GitHub可重算已有派生汇总，并不等于已上传所有原始10x大矩阵、所有scVI权重或E5a/E5b原始预测；涉及原始模型重训或逐spot重新评分的复现仍取决于那些未公开文件的可访问性。

\subsection{研究方法的适用范围}
\label{subsec:cc-method-limits-v1}

本研究把CellCharter作为成熟的空间表征工具，所设计的是特征深度、聚类数、边界精度与外部评价的系统实验，不以新网络、新聚类损失或经证实的算法改良为核心成果。现有设计有若干同时存在的限制：25D配置及2--4层比较与CellCharter原作者的DLPFC调参和基准重叠，且该配置由原作者在3张切片上按标签调出；候选B在开发切片上提名，后又在含这些切片的完整网格中评价；累计邻域深度增加也增加了特征维数；Oracle-K使用参考类别数先验；不同成熟空间算法的原生预处理和计算成本并不一致；BayesSpace与GMM的K-blind似然判据频繁命中搜索上界，其中BayesSpace判据与官方建议不一致；对照集未包含STAGATE；AutoK只重复5次且逐切片运行；DLPFC的独立供体只有3名，且某些旧来源标注/metadata有历史卫生问题；外部乳腺癌的病理区域标签不能直接与皮层分层标签作跨器官绝对精度比较。上述限制不通过后验更换参数或删除负结果解决，而应在结果与讨论中透明呈现。

本章方法阐述对应真实已执行的计算过程、冻结实验方案及标明时间线的事后诊断。若后续需要扩展文献背景、具体软件版本或学校模板要求，须依据实际源文件和正式参考文献补充，不应在没有原文核对的情况下为提升“创新性”而重新命名已有CellCharter操作。
